\title{Controlling Text-to-Image Diffusion by Orthogonal Finetuning}

\begin{abstract}
Text-to-image diffusion models have recently demonstrated exceptional ability in producing realistic images based on textual descriptions. Effectively steering or modulating these sophisticated models for a variety of downstream applications remains a significant unresolved issue. In response, we propose a theoretically sound adaptation technique named Orthogonal Finetuning~(OFT), which allows text-to-image diffusion models to be tailored for specific downstream tasks. In contrast to conventional methods, OFT guarantees preservation of the hyperspherical energy, a metric reflecting the pairwise relationships between neurons constrained to the unit hypersphere. We demonstrate that maintaining this property is pivotal for retaining the semantic generative power inherent to text-to-image diffusion models. To further bolster finetuning robustness, we introduce Constrained Orthogonal Finetuning (COFT), which integrates an explicit constraint on the hypersphere's radius. We focus on two major finetuning scenarios in text-to-image generation: subject-driven generation, aiming to synthesize images of a specified subject in varied contexts using a limited number of input images along with a textual prompt; and controllable generation, which equips the model to incorporate supplementary control inputs. Experimental results validate that our OFT-based approach consistently surpasses existing techniques in both image quality and finetuning efficiency.
\end{abstract}

\section{Introduction}

Advances in text-to-image diffusion models~\cite{saharia2022photorealistic,ramesh2022hierarchical,rombach2022high} have enabled state-of-the-art text-based high-quality image synthesis. However, text prompts alone often fail to deliver the nuanced and precise control required for certain applications, as they may be inherently vague or under-specified. To overcome these limitations, this work addresses two principal tasks in text-to-image generation:

\begin{itemize}[leftmargin=*,nosep]

    \item \textbf{Subject-driven generation}~\cite{ruiz2023dreambooth}: Using a small set of subject images, the aim is to create new images featuring the same subject, recontextualized according to a guiding text description.
    \item \textbf{Controllable generation}~\cite{zhang2023adding,mou2023t2i}: Leveraging auxiliary control information (\emph{e.g.}, edge detections, segmentation masks), this task seeks to direct the generation process so that outputs conform to both the specified control signal and the given text prompt.
\end{itemize}

At their core, both tasks focus on the challenge of finetuning text-to-image diffusion models in a way that maintains the generative strengths acquired during pretraining. An effective finetuning strategy should satisfy two main criteria: (1) \emph{training efficiency}, which involves minimizing the number of trainable parameters and reducing the training epochs required; and (2) \emph{generalizability preservation}, meaning that the model should retain its capacity for high-quality, diverse image generation. Common finetuning techniques involve either slight adjustments to neuron weights using a modest learning rate (\emph{e.g.}, \cite{ruiz2023dreambooth}), or the introduction of small auxiliary modules with re-parameterized weights (\emph{e.g.}, \cite{hulora2022,zhang2023adding}). While these methods are straightforward, neither approach ensures that the original generative capabilities from pretraining are preserved. Moreover, there remains little theoretical guidance for developing robust finetuning methods or for selecting optimal hyperparameters, such as the appropriate number of training epochs. A central challenge is the absence of an established metric to evaluate how well the pretrained generative functionality is maintained. Most existing approaches operate under the implicit belief that minimizing the Euclidean distance between finetuned and pretrained models translates to better retention of generative abilities; accordingly, very low learning rates are employed. However, we contend that solely relying on the Euclidean distance does not adequately reflect the extent of semantic preservation. Incorporating a more structurally-informed metric to assess the difference between the pretrained and finetuned models could significantly enhance both the stability of finetuning and the retention of pretraining performance.

Drawing from empirical findings that hyperspherical similarity serves as a strong proxy for semantic relationships~\cite{liu2017deep,liu2018decoupled,chen2020angular}, we leverage hyperspherical energy~\cite{liu2018learning} to describe the inter-neuronal relational patterns within a layer. Specifically, hyperspherical energy is calculated as the aggregate of hyperspherical similarities (\emph{e.g.}, cosine similarities) across all neuron pairs within a layer, effectively reflecting how uniformly neurons are distributed on the unit hypersphere~\cite{liu2021learning}. Our working assumption is that optimal finetuning should result in minimal changes to the hyperspherical energy relative to the original, pretrained model. A straightforward approach would be to include a regularization term aimed at maintaining constant hyperspherical energy during finetuning. However, this method lacks theoretical assurance for effectively minimizing the actual hyperspherical energy difference. 

Instead, we utilize an invariance property of hyperspherical energy: the pairwise hyperspherical similarities remain invariant when a shared orthogonal transformation is applied to all neurons. This insight leads us to introduce Orthogonal Finetuning~(OFT), a method for adapting large-scale text-to-image diffusion models to new tasks without altering the underlying hyperspherical energy. The core mechanism involves learning a common orthogonal transformation at the layer level, thus maintaining the pairwise angular relationships among neurons. Conceptually, OFT can be interpreted as redefining the canonical basis for the neurons within a given layer. Furthermore, acknowledging that finetuned models closer in Euclidean distance to their pretrained counterparts tend to better retain original performance, we extend our approach with Constrained Orthogonal Finetuning~(COFT), which restricts the adapted model to a fixed-radius hypersphere centered at the pretrained neuron positions.

The rationale behind employing orthogonal transformations for finetuning neurons is inspired in part by concepts from the 2D Fourier transform, where an image is represented in terms of its magnitude and phase spectra. Notably, the phase spectrum—encoding angular relationships between inputs and basis elements—carries most of the image’s semantic content. For instance, combining an image’s phase spectrum with an arbitrary magnitude spectrum still allows the reconstruction of an image with intact semantics. This illustrates that altering neuron orientations can significantly influence the semantic aspects of generated images, which is a central aim in both subject-driven and controllable generation contexts. Nevertheless, modifying neuron orientations without restraint risks undermining the generative capacity imparted during pretraining. To mitigate this risk, we propose maintaining the angular relationships among all neuron pairs, hypothesizing that these angles are integral to the neural network's internalized knowledge. This motivates our strategy of learning a layer-wise orthogonal transformation for each layer's neurons so as to conserve the hyperspherical energy.

Furthermore, our approach is motivated by orthogonal over-parameterized training~\cite{liu2021orthogonal}, which entails training classification neural networks from randomly initialized states using orthogonal transformations. Such initial networks inherently possess low expected hyperspherical energy, and the objective in~\cite{liu2021orthogonal} is to preserve this low energy throughout training—an approach shown to enhance generalization performance~\cite{liu2018learning,lin2020regularizing}. Their findings establish that orthogonal transformation offers sufficient adaptability to train well-generalizing classification networks. Unlike their focus, our work centers on finetuning text-to-image diffusion models to enable finer control and superior generative performance on downstream tasks. It is important to note two main distinctions between OFT and~\cite{liu2021orthogonal}. Firstly, whereas~\cite{liu2021orthogonal} explicitly reduces hyperspherical energy, OFT is designed to maintain the hyperspherical energy established by the pretrained model, thereby safeguarding its foundational structure during finetuning. Specifically, further minimizing hyperspherical energy in diffusion models may disrupt essential semantic representations. Secondly, OFT is constructed to limit divergence from the pretrained weights, a constraint absent from the method in~\cite{liu2021orthogonal}. Finetuning thus hinges on an effective compromise between adaptability and structural stability, and we contend that OFT achieves this balance. Our main contributions are as follows:

\begin{itemize}[leftmargin=*,nosep]

    \item We introduce Orthogonal Finetuning, a new finetuning approach designed to enhance controllability in text-to-image diffusion models. To further ensure robustness, we present a constrained variant that restricts the angular displacement from the original pretrained model.
    \item In contrast to current finetuning techniques, OFT optimizes the model parameters while guaranteeing the conservation of hyperspherical energy—a property we empirically establish as essential for maintaining the pretrained model’s generative semantics.
    \item We evaluate OFT in the contexts of subject-driven generation and controllable generation, undertaking a detailed experimental analysis that reveals substantial gains over previous approaches in generation fidelity, rate of convergence, and stability during finetuning. Additionally, OFT demonstrates improved sample efficiency, successfully converging with fewer training samples and epochs.
    \item For the task of controllable generation, we introduce a novel control consistency metric to quantitatively measure controllability. This metric is computed by inferring the control signal from the produced image and comparing it to the original input control. Our results corroborate the strong controllability afforded by OFT.
\end{itemize}

\section{Related Work}

\textbf{Text-to-image diffusion models}. Recent advances~\cite{saharia2022photorealistic,ramesh2022hierarchical,rombach2022high,gu2022vector,nichol2022glide} have propelled the field of text-to-image generation, largely propelled by progress in diffusion-based generative techniques~\cite{ho2020denoising,song2021score,song2021denoising,dhariwal2021diffusion} and multimodal representation learning~\cite{su2020vlbert,lu2019vilbert,radford2021learning,wang2021simvlm,li2022blip,li2023blip,alayrac2022flamingo,schuhmann2022laion}. GLIDE~\cite{nichol2022glide} and Imagen~\cite{saharia2022photorealistic} build diffusion models operating directly in pixel space. GLIDE couples text encoder training with a diffusion prior, leveraging paired image-text data, while Imagen retains a fixed pretrained text encoder during training. Models such as Stable Diffusion~\cite{rombach2022high} and DALL-E2~\cite{ramesh2022hierarchical} function in the latent space; Stable Diffusion employs VQ-GAN~\cite{esser2021taming} to construct a visual latent codebook, whereas DALL-E2 utilizes CLIP~\cite{radford2021learning} for a joint embedding space unifying images and texts. Beyond diffusion approaches, alternative generative strategies including adversarial frameworks~\cite{reed2016generative,zhang2017stackgan,xu2018attngan,li2019controllable} and autoregressive architectures~\cite{ramesh2021zero,ding2021cogview,wu2022nuwa,yuscaling} have also addressed text-to-image synthesis. Notably, OFT is model-agnostic and can be integrated with any text-to-image diffusion framework.

\textbf{Subject-driven generation}. To limit unintended alterations of subject content, prior works~\cite{avrahami2022blended,nichol2022glide} incorporate user-supplied masks as auxiliary inputs. Inversion-based approaches~\cite{choi2021ilvr,dhariwal2021diffusion,ramesh2022hierarchical,gal2022image} modify image context while leaving the central subject intact. \cite{hertz2022prompt} achieves local and global image editing without the need for user-defined masks. Nevertheless, these strategies often fail to robustly retain subject identity details. Pivotal Tuning~\cite{roich2022pivotal} addresses this by finetuning a generator at the neighborhood of an initial inverted latent code, using a dedicated regularizer to maintain identity fidelity. Likewise, \cite{nitzan2022mystyle} formulates a personalized facial generative prior using multiple images of an individual. \cite{casanova2021instance} permits the synthesis of diverse instance variants, though sometimes at the expense of instance-specific accuracy. DreamBooth~\cite{ruiz2023dreambooth} leverages both a custom textual token and a small set of subject images, applying a reconstruction loss alongside a class-specific prior loss for training the diffusion model. In our approach, OFT follows the DreamBooth setup but replaces conventional small-step parameter updates with orthogonal transformations during finetuning.

\textbf{Controllable generation}. Image-to-image translation falls under the broader umbrella of controllable generation. Traditionally, most approaches utilize conditional generative adversarial networks for this purpose~\cite{isola2017image,zhu2017toward,wang2018high,park2019semantic,choi2018stargan}. In addition, diffusion models have emerged as promising alternatives for this task~\cite{saharia2022palette,voynov2022sketch,wang2022pretraining}. Notably, ControlNet~\cite{zhang2023adding} introduces a technique to steer a pretrained diffusion model by finetuning it to accommodate supplementary control inputs, thereby achieving notable performance in controllable generation. Similarly, the contemporaneous work T2I-Adapter~\cite{mou2023t2i} refines pretrained diffusion models to boost the controllability of output images. Building on the task setup described in \cite{zhang2023adding,mou2023t2i}, we leverage OFT to finetune pretrained diffusion models, attaining superior controllability with reduced training data and a lower parameter count for finetuning. Importantly, OFT does not impose any extra computation cost during inference.

\textbf{Model finetuning}. Adapting large, pretrained models to specific downstream tasks via finetuning has become widespread in recent years~\cite{devlin2018bert,brown2020language,he2022masked}. Within the realm of finetuning, numerous adaptation strategies (\emph{e.g.}, \cite{hulora2022,houlsby2019parameter,pfeiffer2020adapterfusion}) have been intensively explored, particularly in natural language processing. Among these, LoRA~\cite{hulora2022} stands out in relation to OFT, employing a low-rank decomposition for additive weight updates during finetuning. In contrast, OFT applies a layer-wise, shared orthogonal transformation to modify neuron weights multiplicatively; this approach mathematically guarantees the preservation of pairwise angles between neurons within the same layer, resulting in increased stability.

\section{Orthogonal Finetuning}

\subsection{Why Does Orthogonal Transformation Make Sense?}

We first examine the motivation for using orthogonal transformations in the finetuning of text-to-image diffusion models. To clarify, this issue is divided into two aspects: (1) the rationale for adjusting neuron angles (\emph{i.e.}, their directions) during finetuning, and (2) the justification for specifically employing orthogonal transformations for modifying these angles.

To address the first question, we take cues from empirical findings presented in \cite{liu2018decoupled,chen2020angular}, which suggest that the angular differences between features effectively capture the semantic distinction. SphereNet~\cite{liu2017deep} further demonstrates that when all neurons in a neural network are normalized to reside on a unit hypersphere, the network retains similar representational capacity and often achieves enhanced generalization. This observation highlights that the orientation of neurons encodes critical information from the data. To illustrate the significance of neuron orientation, we design a simple experiment as shown in Figure~\ref{fig:toy_ae}, where a conventional convolutional autoencoder is trained on a dataset of flower images. During training, the autoencoder constructs feature maps via the typical inner product (with $z$ representing the convolution between kernel $\bm{w}$ and input patch $\bm{x}$). For evaluation, we compare three feature extraction schemes: (a) the training-phase inner product, (b) preserving only the feature magnitude, and (c) retaining solely the angular component. According to Figure~\ref{fig:toy_ae}, reconstructing images using angular features alone achieves near-perfect input recovery, whereas retaining only the feature magnitude fails to yield meaningful information. Notably, the network is trained exclusively with the inner product—no cosine-based or angular activation is applied in training—emphasizing that neuron directions are primarily responsible for conveying the input’s semantic content. Thus, to manipulate semantic information embedded in the images, altering neuron orientations may offer a particularly effective strategy.

Turning to the second question, the most direct approach for adjusting neuron orientation involves applying a common rotation or reflection to all neurons within a layer, inherently invoking orthogonal transformations. Although alternative angular manipulations—such as rotating individual neurons by distinct angles—could be considered for greater adaptability, our findings indicate that orthogonal transformations present an optimal balance between structural regularity and expressive flexibility. Furthermore, as demonstrated in \cite{liu2021orthogonal}, orthogonal transformations possess sufficient expressive power for neural network optimization. To reinforce this point, we conduct a regularization experiment using the ControlNet~\cite{zhang2023adding} controllable generation framework; corresponding results are depicted in Figure~\ref{fig:ortho}. It can be observed that when orthogonal transformations are imposed, our standard OFT method remains stable and achieves precise generation control by epoch 20. In contrast, omitting the orthogonality constraint in OFT results in an inability to generate realistic samples or maintain controllability. These experimental outcomes underscore the essential role of orthogonality in effective OFT.

In standard finetuning approaches, gradient descent is generally employed with a reduced learning rate to update either the entire model or selected layers. Utilizing a small learning rate implicitly favors minimal departures from the pretrained parameters, such that conventional finetuning effectively conducts optimization under an implicit penalty term $\thickmuskip=2mu \medmuskip=2mu \| \bm{M} - \bm{M}^0\|$, where $\bm{M}$ represents the adapted weights and $\bm{M}^0$ denotes the pretrained parameters. While such implicit regularization has an intuitive appeal, it may nonetheless grant excessive flexibility when dealing with large-scale models. To mitigate this, LoRA adds an explicit low-rank restriction on weight modifications, namely $\thickmuskip=2mu \medmuskip=2mu \text{rank}(\bm{M} - \bm{M}^0)=r'$, setting $r'$ to a small predefined value. In contrast, OFT imposes a constraint focused on pairwise neuron similarities, requiring $\thickmuskip=2mu \medmuskip=2mu \| \text{HE}(\bm{M})-\text{HE}(\bm{M}^0) \|=0$, where $\text{HE}(\cdot)$ represents the hyperspherical energy for a given weight matrix.

Consider, for illustration, a fully connected layer characterized by weight matrix $\thickmuskip=2mu \medmuskip=2mu \bm{W}=\{\bm{w}_1,\cdots,\bm{w}_n\}\in\mathbb{R}^{d\times n}$ with neurons $\bm{w}_i\in\mathbb{R}^{d}$ and pretrained version $\bm{W}^0$. The output $\thickmuskip=2mu \medmuskip=2mu \bm{z}\in\mathbb{R}^n$ for input $\thickmuskip=2mu \medmuskip=2mu \bm{x}\in\mathbb{R}^d$ is $\bm{z} = \bm{W}^\top\bm{x}$. The OFT framework can thus be interpreted as driving the difference in hyperspherical energy between the finetuned and original models toward zero:

\begin{equation}\label{eq:oft_principle}
\footnotesize
\min_{\bm{W}} \left\| \text{HE}(\bm{W})-\text{HE}(\bm{W}^0)\right\|~~\Leftrightarrow~~\min_{\bm{W}} \bigg{\|} \sum_{i\neq j} \|\hat{\bm{w}}_i-\hat{\bm{w}}_j\|^{-1}-\sum_{i\neq j} \|\hat{\bm{w}}^0_i-\hat{\bm{w}}^0_j\|^{-1}\bigg{\|}
\end{equation}

with each normalized neuron defined as $\thickmuskip=2mu \medmuskip=2mu \hat{\bm{w}}_i:=\bm{w}_i/\|\bm{w}_i\|$, and hyperspherical energy for the weight matrix $\bm{W}$ given by $\thickmuskip=2mu \medmuskip=2mu \text{HE}(\bm{W}):= \sum_{i\neq j} \|\hat{\bm{w}}_i-\hat{\bm{w}}_j\|^{-1}$. The minimal possible value of Eq.~\eqref{eq:oft_principle} is evidently zero, which is attainable if $\bm{W}$ and $\bm{W}^0$ only vary by a rotation or reflection, i.e., $\thickmuskip=2mu \medmuskip=2mu \bm{W} = \bm{R}\bm{W}^0$, where $\thickmuskip=2mu \medmuskip=2mu \bm{R}\in\mathbb{R}^{d\times d}$ is an orthogonal matrix ($\det \bm{R}=1$ denotes rotation, while $\det \bm{R}=-1$ indicates reflection). The central premise of OFT, then, is to adapt the network by optimizing a shared orthogonal matrix per layer, which transforms the neurons within each layer. In formal terms, for a pretrained fully connected layer $\thickmuskip=2mu \medmuskip=2mu\bm{W}^0\in\mathbb

\begin{equation}\label{eq:oft_general}
    \footnotesize
    \bm{z}=\bm{W}^\top\bm{x}=(\bm{R}\cdot\bm{W}^0)^\top\bm{x},~~~\text{s.t.}~\bm{R}^\top\bm{R}=\bm{R}\bm{R}^\top=\bm{I}
\end{equation}

Here, $\bm{W}$ represents the weight matrix adapted through OFT-finetuning, while $\bm{I}$ signifies the identity matrix. An overview of OFT is provided in Figure~\ref{fig:block}. Following the practice of zero initialization used in LoRA, it is crucial that OFT modifies the pretrained model strictly starting from $\bm{W}^0$. This requirement is fulfilled by initializing the orthogonal matrix $\bm{R}$ as the identity matrix, thereby ensuring the finetuned model begins with the original pretrained weights. To maintain the orthogonality property of $\bm{R}$, we employ differentiable orthogonalization methods referenced in \cite{liu2021orthogonal,lezcano2019cheap}. We later elaborate on how orthogonality is preserved efficiently in computation.

\subsection{Efficient Orthogonal Parameterization}\label{sect:eff_ortho}

Conventional orthogonalization techniques, such as the Gram-Schmidt process, though differentiable, tend to be computationally prohibitive for large-scale models~\cite{liu2021orthogonal}. To address computational bottlenecks, we make use of the Cayley parameterization for constructing orthogonal matrices. Concretely, the orthogonal matrix is parameterized as $\thickmuskip=2mu \medmuskip=2mu \bm{R}=(\bm{I}+\bm{Q})(I-\bm{Q})^{-1}$, where the matrix $\bm{Q}$ is skew-symmetric, i.e., $\thickmuskip=2mu \medmuskip=2mu \bm{Q}=-\bm{Q}^\top$. The Cayley parameterization introduces an efficiency gain at the minor expense that it can only yield orthogonal matrices with determinant equal to $1$, restricting $\bm{R}$ to the special orthogonal group. In practical scenarios, we observe that this restriction has negligible impact on empirical results. However, even with Cayley parameterization, representing $\bm{R}$ directly can become storage-intensive as $d$ increases. To mitigate this, we introduce a block-diagonal structure to $\bm{R}$, dividing it into $r$ blocks, as formulated below:

\begin{equation}
    \footnotesize
    \bm{R}=\text{diag}(\bm{R}_1,\bm{R}_2,\cdots,\bm{R}_r)=\begin{bmatrix}
    \bm{R}_1\in \text{O}(\frac{d}{r}) &  &\\
    &\ddots & \\
    & & \bm{R}_r\in \text{O}(\frac{d}{r})
    \end{bmatrix}\in \text{O}(d)
\end{equation}

Here, $\text{O}(d)$ refers to the orthogonal group of dimension $d$, where $\thickmuskip=2mu \medmuskip=2mu \bm{R}\in\mathbb{R}^{d\times d}$ and, for every $i$, $\thickmuskip=2mu \medmuskip=2mu \bm{R}_i\in\mathbb{R}^{d/r\times d/r}$ are orthogonal matrices. When $\thickmuskip=2mu \medmuskip=2mu r=1$, the block-diagonal structure collapses to a full, unconstrained orthogonal matrix. A $d\times d$ orthogonal matrix possesses $\thickmuskip=2mu \medmuskip=2mu d(d-1)/2$ independent parameters, resulting in $\mathcal{O}(d^2)$ computational complexity. In contrast, for an orthogonal matrix structured as $r$-block diagonal, the parameter count becomes $\thickmuskip=2mu \medmuskip=2mu d(d/r-1)/2$ with corresponding complexity $\thickmuskip=2mu \medmuskip=2mu \mathcal{O}(d^2/r)$. To further economize on parameter usage, one may optionally tie the block matrices together (\emph{i.e.}, $\thickmuskip=2mu \medmuskip=2mu\bm{R}_i=\bm{R}_j$ for all $i\neq j$), thereby reducing the parameter complexity to $\mathcal{O}(d^2/r^2)$. Importantly, despite these efficiency measures, the overall orthogonality of $\bm{R}$ is preserved, which maintains hyperspherical energy conservation.

We now compare the parameter efficiency of OFT with that of LoRA. LoRA, characterized by a low-rank parameter $r'$, incurs a trainable parameter count of $\thickmuskip=2mu \medmuskip=2mu r'(d+n)$. If both $r$ and $r'$ are assumed to scale with the neuron dimensionality $d$ (\emph{e.g.}, $\thickmuskip=2mu \medmuskip=2mu r=r'=\alpha d$ for some fixed constant $\thickmuskip=2mu \medmuskip=2mu 0<\alpha\leq 1$), then LoRA’s parameter complexity is $\mathcal{O}(d^2+dn)$, while OFT’s is reduced to $\mathcal{O}(d)$. To concretely demonstrate this disparity, consider a weight matrix of size $\thickmuskip=2mu \medmuskip=2mu 128\times 128$: with $\thickmuskip=2mu \medmuskip=2mu r'=8$, LoRA trains $2,048$ parameters, whereas OFT (with $r=8$ and without block sharing) only requires $960$ trainable parameters.

\subsection{Constrained Orthogonal Finetuning}

It is possible to further restrict the expressiveness of the OFT approach by imposing a requirement that the finetuned parameters reside close to those of the original pretrained model. COFT implements this by utilizing the following forward computation:

\begin{equation}\label{eq:coft_general}
    \footnotesize
    \bm{z}=\bm{W}^\top\bm{x}=(\bm{R}\cdot\bm{W}^0)^\top\bm{x},~~~\text{s.t.}~\bm{R}^\top\bm{R}=\bm{R}\bm{R}^\top=\bm{I},~\norm{\bm{R}-\bm{I}}\leq \epsilon
\end{equation}

Here, $\bm{R}$ must satisfy both an orthogonality requirement and a constraint that bounds its deviation from the identity matrix by $\epsilon$. While the Cayley parameterization presented in Section~\ref{sect:eff_ortho} can be employed to ensure orthogonality, directly integrating the $\epsilon$-deviation restriction into a Cayley-parameterized $\bm{R}$ is not straightforward. To investigate this, we utilize the Neumann series for approximating $\thickmuskip=2mu \medmuskip=2mu \bm{R}=(\bm{I}+\bm{Q})(I-\bm{Q})^{-1}$, resulting in the approximation $\thickmuskip=2mu \medmuskip=2mu \bm

Since the series converges under the operator norm, it is permissible to incorporate the constraint $\thickmuskip=2mu \medmuskip=2mu\|\bm{R}-\bm{I}\|\leq\epsilon$ directly within the Cayley transform. Under this transformation, the corresponding constraint becomes $\thickmuskip=2mu \medmuskip=2mu \|\bm{Q}-\bm{0}\|\leq\epsilon'$, where $\bm{0}$ stands for a matrix of all zeros and $\epsilon'$ denotes a distinct error threshold, not necessarily equal to $\epsilon$. This updated constraint on $\bm{Q}$ lends itself naturally to enforcement via projected gradient descent. To ensure that the orthogonal matrix $\bm{R}$ is initially the identity, $\bm{Q}$ is set to the zero matrix at initialization. Conceptually, COFT incorporates two explicit forms of regularization: it maintains a low difference in hyperspherical energy and restricts the deviation from the pretrained weights. While the latter is typically an implicit assumption in most finetuning strategies, COFT formalizes it as a concrete constraint. Our empirical results indicate that, although OFT already provides robust performance, introducing this explicit deviation regularization in COFT further enhances stability during finetuning. As demonstrated in Figure~\ref{fig:eps_coft}, varying $\epsilon$ adjusts the degree of permissible deviation: a higher $\epsilon$ yields COFT outcomes close to those of OFT, whereas reducing $\epsilon$ causes the COFT-finetuned model to act more like the original pretrained text-to-image diffusion model.

\subsection{Re-scaled Orthogonal Finetuning}

We introduce an augmentation to the standard OFT framework by jointly optimizing a neuron-wise scaling factor in addition to the orthogonal transformation. The rationale is that scaling neuron outputs leaves hyperspherical energy invariant, as magnitudes are normalized during evaluation. Concretely, we reformulate the forward computation as $\thickmuskip=2mu \medmuskip=2mu\bm{z}=(\bm{R}\bm{W}^0\bm{D})^\top\bm{x}$\footnote{\textbf{Errata}: In the NeurIPS camera ready submission, we mistakenly described the forward computation for the re-scaled OFT as $\thickmuskip=2mu \medmuskip=2mu\bm{z}=(\bm{D}\bm{R}\bm{W}^0)^\top\bm{x}$. The original code implementation remains correct; thus, the findings in Appendix~\ref{app:rescaled} are valid.} with $\thickmuskip=2mu \medmuskip=2mu \bm{D}=\text{diag}(s_1,\cdots,s_n)\in\mathbb{R}^{n\times n}$, where $\bm{D}$ is a diagonal matrix whose diagonal entries $s_1,\cdots,s_n$ are all positive and learnable. Unlike the original OFT update in Eq.~\eqref{eq:oft_general}, which only adjusts $\bm{R}$, both $\bm{R}$ and $\bm{D}$ are learnable parameters in this re-scaled variant. This extension slightly increases parameter count, but further enhances OFT's capacity. Although our experiments primarily rely on the unmodified OFT to clearly isolate the impact of orthogonal updates, we consistently observe that the re-scaled version delivers even better performance (refer to Appendix~\ref{app:rescaled}).

\section{Intriguing Insights and Discussions}

\textbf{OFT's applicability is model-agnostic}. Fundamentally, OFT is compatible with a wide range of neural architectures. For instance, whereas LoRA is commonly tailored for the attention layers of Transformers~\cite{hulora2022}, our evaluation targets only the attention weights in order to ensure comparability. In addition to fully connected layers, OFT is naturally compatible with convolutional modules. Unlike LoRA, the block-diagonal form of $\bm{R}$ has intuitive interpretations within convolutional architectures. If we construct as many blocks as there are input channels, each block targets a specific neuron channel, mirroring the design of depth-wise convolution filters~\cite{chollet2017xception}. Conversely, sharing all $\bm{R}$ blocks across channels applies a single orthogonal transformation to all neurons. We provide an initial exploration of applying OFT to convolutional layers in Appendix~\ref{app:convolution}.

\textbf{Connection to LoRA}. LoRA introduces a low-rank matrix to the pretrained weight matrix to avoid significant alterations to the existing learned parameters. In contrast, OFT constrains the transformation applied to the pretrained weight matrix by enforcing it to be orthogonal (and therefore full-rank), thereby safeguarding the integrity of pretrained knowledge. The OFT forward computation can be reformulated as $\thickmuskip=2mu \medmuskip=2mu \bm{z}=(\bm{R}\bm{W}^0)^\top\bm{x}=(\bm{W}^0+(\bm{R}-\bm{I})\bm{W}^0)^\top\bm{x}$, where $\thickmuskip=2mu \medmuskip=2mu (\bm{R}-\bm{I})\bm{W}^0$ resembles the low-rank adaptation employed by LoRA. Given that $\bm{W}^0$ is usually of full rank, OFT also induces a low-rank update when $\thickmuskip=2mu \medmuskip=2mu \bm{R}-\bm{I}$ is itself of low rank. LoRA introduces a tunable rank parameter $r'$ to control parameter efficiency, while OFT adopts a block-diagonal structure with parameter $r$ to similarly reduce the quantity of parameters requiring optimization. Importantly, LoRA and OFT offer fundamentally different parameter-efficient adaptation mechanisms: LoRA leverages low-rank decompositions to economize on trainable parameters, whereas OFT exploits the sparsity pattern of block-diagonal orthogonal matrices.

\textbf{Why OFT converges faster}. On the one hand, Figure~\ref{fig:toy_ae} demonstrates that manipulating the directions of neurons is the most direct strategy for modifying semantic content, a goal that OFT targets explicitly. On the other hand, OFT's parameterization allows fine-tuning on the manifold of a smooth hypersphere, facilitating a more tractable and favorable optimization landscape. Such behavior has also been confirmed through empirical observation in \cite{liu2021orthogonal}.

\textbf{Why not minimize hyperspherical energy}. Our approach differs from \cite{liu2021orthogonal} in that we do not target hyperspherical energy minimization. For classification tasks, the absence of redundant neurons is beneficial, but enforcing a configuration of minimal hyperspherical energy, which uniformly distributes neurons over the hypersphere, may be detrimental to the retention of pretrained knowledge. Thus, such an objective is not suitable for finetuning purposes.

\textbf{Balancing flexibility and regularity in finetuning}. Our analysis reveals an inherent trade-off between the flexibility and regularity of finetuning approaches. While standard finetuning offers maximal flexibility, it often comes at the expense of model stability and is prone to collapse. In contrast, OFT, despite its simplicity, achieves a favorable compromise by maintaining the pairwise angular relationships between neurons. Additionally, the block-diagonal approach can be interpreted as a more stringent form of regularization for the orthogonal matrix.

\textbf{Zero inference overhead}. Differing from ControlNet, the OFT framework imposes no additional inference burden on the finetuned network. During inference, the learned orthogonal transformation matrix $\bm{R}$ can be directly multiplied with the pretrained weight matrix $\bm{W}^0$ to yield the effective weight matrix $\thickmuskip=2mu \medmuskip=2mu \bm{W}=\bm{R}\bm{W}^0$. This ensures that inference efficiency matches that of the original pretrained model.

\section{Experiments and Results}

\textbf{Experimental setup}. For our empirical study, Stable Diffusion v1.5~\cite{rombach2022high} serves as the backbone pretrained text-to-image system. In order to maintain fairness, we randomly select images generated from each technique. The experimental protocol for subject-driven generation aligns closely with DreamBooth~\cite{ruiz2023dreambooth}, whereas controllable generation settings are based on ControlNet~\cite{zhang2023adding} and T2I-Adapter~\cite{mou2023t2i}. To provide a balanced comparison against LoRA, OFT and COFT are restricted to operate on the same network layers as those adapted by LoRA. Additional implementation specifics and results are included in Appendix~\ref{app:settings}.

\subsection{Subject-driven Generation}

\textbf{Configuration}. The experiments use DreamBooth~\cite{ruiz2023dreambooth} and LoRA~\cite{hulora2022} as baseline models. All methods share the identical loss objective defined by DreamBooth. Where relevant, we adhere to each method’s official guidelines and optimal hyperparameter configurations. Supplementary outcomes and implementation details can be found in Appendix~\ref{app:settings},\ref{app:coft_oft},\ref{app:more_qualitative},\ref{app:failure}.

\textbf{Finetuning stability and convergence}. We begin by assessing the robustness and convergence rate during finetuning for DreamBooth, LoRA, OFT, and COFT. The findings are illustrated in Figure~\ref{fig:teaser} and Figure~\ref{fig:db_convergence}. It is evident from these results that COFT and OFT both provide stable finetuning for the diffusion model. Within 400 iterations, both DreamBooth and OFT-based methods manage effective control over the generation process, whereas LoRA demonstrates a deficiency in maintaining the subject’s identity. Upon extending to 2000 iterations, DreamBooth exhibits mode collapse in its outputs, and LoRA misattributes features, producing yellow fur instead of a yellow shirt. In sharp contrast, OFT and COFT continue to deliver reliable and controlled image generation throughout this process. This supports the claim that OFT ensures not only rapid but also dependable convergence in subject-driven generation scenarios. Importantly, the robustness with respect to the total number of finetuning iterations is highly beneficial, as it simplifies the practical challenge of tuning iteration counts for various subjects. With both OFT and COFT, a generously large iteration count can be preset without the need for intricate calibration. In the case of COFT, once an appropriate $\epsilon$ is chosen, both the learning rate and the iteration count require minimal manual adjustment.

\textbf{Quantitative comparison}. Adopting the evaluation protocol in \cite{ruiz2023dreambooth}, we quantitatively assess our methods using subject fidelity metrics (DINO~\cite{caron2021emerging}, CLIP-I~\cite{radford2021learning}), text prompt alignment (CLIP-T~\cite{radford2021learning}), and sample diversity (LPIPS~\cite{zhang2018unreasonable}). The CLIP-I metric represents the mean pairwise cosine similarity between CLIP embeddings of generated and actual images. For DINO, a similar computation is performed using ViT S/16 DINO embeddings. The CLIP-T metric involves averaging the cosine similarity between text prompt and generated image embeddings. To evaluate diversity, we consider the average LPIPS cosine distance among generated samples for a single subject under identical textual conditions. As indicated by Table~\ref{table:dreambooth_final}, COFT and OFT significantly outperform DreamBooth and LoRA on DINO and CLIP-I measures, with comparable or slightly improved outcomes on prompt fidelity and diversity criteria. Each approach undergoes finetuning from the same pretrained checkpoint with 30 different random seeds to reduce the impact of stochastic variations. These outcomes confirm that the OFT framework attains superior convergence properties and stability, consistently delivering enhanced final results.

\textbf{Qualitative comparison}. To offer a clearer insight into the advantages of OFT, we present a selection of randomly chosen samples for subject-driven generation in Figure~\ref{fig:dreambooth_final}. To ensure fairness, each sample is generated from the same finetuned model under different methods, with no separate prompt optimization for any approach. The best-performing model for each method, as measured by the highest validation CLIP score, is used. The examples in Figure~\ref{fig:dreambooth_final} indicate that both OFT and COFT maintain strong semantic fidelity to the subject, whereas LoRA frequently fails to capture the subject’s identity (for instance, in the bowl scenario, LoRA does not preserve the original subject at all). Furthermore, OFT and COFT display a more precise response to text prompts, unlike DreamBooth, which, though effective at retaining subject characteristics, often neglects the input prompt (as demonstrated in the bowl case, first row). These qualitative results collectively show that OFT achieves superior balance between subject fidelity and prompt-based controllability. In addition, OFT is robust to the number of iterations, showing consistently high performance even as iterations increase—a property not shared by DreamBooth or LoRA. Additional examples can be found in Appendix~\ref{app:more_qualitative}. We also conduct human evaluations in Appendix~\ref{app:human}, which further corroborate the superiority of OFT.

\subsection{Controllable Generation}

\textbf{Settings}. For baseline comparison, we adopt ControlNet~\cite{zhang2023adding}, T2I-Adapter~\cite{mou2023t2i}, and LoRA~\cite{hulora2022}. Our evaluation spans three difficult controllable generation scenarios: canny edge to image~(C2I) on the COCO dataset~\cite{lin2014microsoft}, segmentation map to image~(S2I) using ADE20K~\cite{zhou2017scene}, and landmark to face~(L2F) on CelebA-HQ~\cite{karras2018progressive,xia2021tedigan}. All techniques are applied to finetune Stable Diffusion~(SD) v1.5 on these datasets for 20 training epochs. Expanded qualitative results are provided in Appendix~\ref{app:more_qualitative},~\ref{app:more_control_task}, and~\ref{app:failure}.

\textbf{Convergence}. We assess how rapidly ControlNet, T2I-Adapter, LoRA, and COFT converge when applied to the L2F task, employing both quantitative and qualitative methods. For the quantitative assessment, we use the mean $\ell_2$ distance between the control and predicted face landmarks as our metric. As illustrated in Figure~\ref{fig:face_convergence}, we present the progression of the face landmark error as each model is finetuned across varying numbers of epochs. It is evident from the results that both COFT and OFT attain considerably faster convergence relative to the baselines. Notably, LoRA requires 20 epochs to achieve the same level of performance that our OFT framework reaches after just 8 epochs. It is important to highlight that OFT and COFT employ a number of trainable parameters on par with LoRA (which itself is substantially fewer than ControlNet), while converging significantly more efficiently than previous approaches. Furthermore, the rapid convergence capability of OFT is corroborated by the results displayed in Figure~\ref{fig:teaser}. The rightmost example in this figure demonstrates the superior data-efficiency of OFT compared to both ControlNet and LoRA, as OFT reaches effective convergence using only 5\% of the ADE20K dataset. On the qualitative side, we primarily contrast OFT, COFT, and ControlNet, since the latter exhibits landmark errors closest to those of our method. The outcomes visualized in Figure~\ref{fig:face_convergence_qual} reveal that both OFT and COFT not only converge reliably, but also yield facial poses that become increasingly aligned with the control landmarks throughout training. In contrast, ControlNet's controllability emerges abruptly post the 8-th epoch, mirroring the abrupt convergence behavior previously reported in \cite{zhang2023adding}. The consistent and smooth convergence offered by our OFT framework greatly enhances practical usability, as its training trajectory is much more stable and foreseeable, making hyperparameter selection considerably more straightforward.

\textbf{Quantitative comparison}. We propose a metric for assessing control consistency to measure how well controllable generation models adhere to input control signals. This metric involves extracting the control cue from the generated output and contrasting it with the original control signal provided as input. Specifically, for the C2I task, we evaluate performance using the IoU and F1 score; for S2I, we report mean IoU along with mean and overall accuracy; and for L2F, we use the mean $\ell_2$ distance between the original control landmarks and those predicted by the model. Additional details about the calculation of these metrics are provided in Appendix~\ref{app:settings}. For all evaluated methods, we adopt optimal hyperparameter configurations. Table~\ref{table:control_gen} presents our findings, revealing that OFT and COFT deliver significantly stronger and more precise control than alternative approaches. We notice that methods relying on adapters (such as T2I-Adapter and ControlNet) demonstrate slower convergence rates and inferior final outcomes. LoRA outperforms ControlNet for the S2I task but trails in the C2I and L2F tasks. Overall, LoRA’s maximal achievable performance appears restricted, regardless of extensive hyperparameter optimization. By contrast, OFT’s results show no signs of plateauing, as its performance continues to improve with increasing numbers of trainable parameters. Notably, our systematic approach to quantitatively evaluating controllable generation is among the key contributions of this work, enabling accurate assessment of model controllability, within the reliability limits of existing segmentation or detection tools.

\textbf{Qualitative comparison}. We further perform a qualitative evaluation of OFT and COFT in relation to leading approaches such as ControlNet, T2I-Adapter, and LoRA. As illustrated by randomly sampled results in Figure~\ref{fig:control_final}, both OFT and COFT generate images that are not only highly realistic and visually detailed, but also closely follow the intended controls. For the S2I task, LoRA fails to align its outputs with the supplied segmentation map, whereas ControlNet, OFT, and COFT successfully guide the generation process. OFT and COFT stand out in comparison to ControlNet by producing more detailed, high-quality, and geometrically plausible images with a reduced number of model parameters. When tackling the C2I task, OFT and COFT effectively synthesize semantically coherent images from rough Canny edge inputs, while T2I-Adapter and LoRA lag behind in performance. Regarding the L2F task, our approach delivers the highest accuracy in pose-conditioned face generation, even under difficult facial orientations. Across all three control scenarios, OFT and COFT consistently outperform established methods in qualitative image quality, underscoring the strength of our OFT architecture for controllable image synthesis. To enable a broader qualitative assessment, we present additional visual samples for the three primary control tasks in Appendix~\ref{app:control_exp}, and further demonstrate the versatility of OFT in handling various control problems—including dense pose to human body, sketch to image, and depth to image—in Appendix~\ref{app:more_control_task}.

\section{Concluding Remarks and Open Problems}

Recognizing that angular relationships between neurons play a pivotal role in shaping visual semantics, this work introduces a straightforward yet powerful finetuning approach—orthogonal finetuning—for modulating text-to-image diffusion models. Our investigation centers on two primary applications: subject-driven generation and controllable generation. When evaluated against prevailing techniques, OFT achieves superior controllability and enhanced finetuning stability, while requiring fewer trainable parameters. Notably, OFT does not add computational complexity during inference, thereby yielding a model suitable for practical deployment.

Several intriguing open questions arise from OFT. The method enforces orthogonality by employing Cayley parametrization, necessitating matrix inversion, which poses some scalability constraints. While our adoption of block diagonal parametrization partially addresses this, the quest for a more efficient, differentiable matrix inversion method remains unresolved. Additionally, OFT exhibits a distinctive promise for compositionality: orthogonal matrices derived from several OFT finetuning tasks may be composed through multiplication to yield another orthogonal matrix. An open area of research concerns whether these orthogonal transformations retain knowledge acquired from each corresponding downstream task. Lastly, the degree of parameter efficiency in OFT largely stems from the use of block diagonal structures, which unavoidably inject extra biases and constrain flexibility. Advancing toward more effective and unbiased strategies for boosting parameter efficiency is an essential topic for future work.

\end{document}